\documentclass[10pt]{article}

\usepackage[T1]{fontenc}
\usepackage[hmargin=2.5cm,vmargin=2.0cm]{geometry}
\usepackage{amsmath,amssymb}
\usepackage{booktabs}
\usepackage{microtype}
\usepackage{mathpazo}
\usepackage[colorlinks=true,linkcolor=blue,urlcolor=blue]{hyperref}

\title{Why Permanent Spending Increases Produced Persistent Deflation}
\author{Spending-efficiency model --- internal investigation note}
\date{July 2026}

\newcommand{\ssval}[1]{{#1}^{SS}}

\begin{document}
\maketitle

\noindent
\textbf{Finding.}
The persistent deflation previously generated by permanent spending increases was a
consequence of the monetary-policy rule, not a direct implication of the model's
productive-spending channels.
The rule responded to output relative to its original steady state.
A permanent increase in equilibrium output therefore created a permanent policy wedge.
The bond Euler equation and the Taylor rule could reconcile that wedge only through
inflation below target.
Setting the output coefficient to zero removes this mechanism and returns inflation to
target in the long run.

\section{The permanent-shock experiment}

The experiment raises one public-spending instrument by 1 percent of GDP for 1,000
quarters.
The increase is therefore effectively permanent over the reported 25-year horizon.
It is debt financed over the same interval, so public debt itself does not reach a new
steady state.
To isolate the nominal mechanism, consider the limiting case in which the real
allocation becomes constant while output remains above its original steady-state level.
This endpoint argument closely approximates the long-horizon responses.

The old Taylor rule was
\begin{equation}\label{eq:taylor-old}
    \frac{R_t}{\ssval{R}}
    =
    \left(\frac{R_{t-1}}{\ssval{R}}\right)^{\rho_R}
    \left[
        \left(\frac{\Pi_t}{\ssval{\Pi}}\right)^{\gamma_\pi}
        \left(\frac{Y_t^d}{\ssval{Y^d}}\right)^{\gamma_y}
    \right]^{1-\rho_R}.
\end{equation}
The output term is a deviation from the \emph{original level} of output.
It is not an output gap measured relative to natural or potential output.
Consequently, it does not close when a permanent supply-side improvement raises the
economy's sustainable output level.

\section{Why the old rule forced deflation}

Define log deviations from the original steady state as
\[
    \widehat r \equiv \log(R/\ssval{R}), \qquad
    \widehat\pi \equiv \log(\Pi/\ssval{\Pi}), \qquad
    \widehat y \equiv \log(Y^d/\ssval{Y^d}).
\]
At a constant endpoint, interest-rate smoothing drops out of
equation~\eqref{eq:taylor-old}.
With no monetary-policy shock, the rule becomes
\begin{equation}\label{eq:taylor-endpoint}
    \widehat r = \gamma_\pi\widehat\pi+\gamma_y\widehat y.
\end{equation}

The stationary bond Euler equation is
\begin{equation}\label{eq:euler}
    \lambda_t
    = \beta\,\frac{\lambda_{t+1}}{g}\frac{R_t}{\Pi_{t+1}}.
\end{equation}
Marginal utility is constant at the new endpoint.
Since $\ssval{R}=g\ssval{\Pi}/\beta$, equation~\eqref{eq:euler} implies
\begin{equation}\label{eq:fisher-endpoint}
    \widehat r=\widehat\pi.
\end{equation}
Combining equations~\eqref{eq:taylor-endpoint} and
\eqref{eq:fisher-endpoint} gives the key restriction:
\begin{equation}\label{eq:puzzle}
    \boxed{
    \widehat\pi
    =-\frac{\gamma_y}{\gamma_\pi-1}\widehat y
    }.
\end{equation}

The old calibration used $\gamma_y=0.25$ and $\gamma_\pi=1.50$.
It therefore imposed $\widehat\pi=-0.5\widehat y$ at the endpoint.
Because the figures annualize quarterly inflation, their inflation response is
approximately minus twice the output response when both are expressed in percent.
This accounts for the apparently large inflation response despite small movements in
current real marginal cost.
The long-run mapping was imposed jointly by monetary policy and the Euler equation;
it was not generated by current marginal cost alone.

\begin{table}[ht]
    \centering
    \small
    \caption{Responses after 25 years under the old and new policy rules}
    \label{tab:comparison}
    \begin{tabular}{lrrrr}
        \toprule
        & \multicolumn{2}{c}{Output (percent)}
        & \multicolumn{2}{c}{Inflation (annualized pp.)} \\
        \cmidrule(lr){2-3}\cmidrule(lr){4-5}
        Instrument & Old rule & New rule & Old rule & New rule \\
        \midrule
        Government consumption    & 0.57 & 0.58 & -1.13 & 0.02 \\
        Infrastructure investment & 2.44 & 2.62 & -4.79 & 0.14 \\
        Human-capital investment  & 4.04 & 4.31 & -7.90 & 0.57 \\
        Public R\&D                & 4.13 & 4.59 & -8.00 & 0.29 \\
        \bottomrule
    \end{tabular}

    \vspace{0.4em}
    \parbox{0.92\textwidth}{\footnotesize
    \textit{Note:} Output is the percent deviation from the original steady state.
    Inflation is the annualized percentage-point deviation.
    The old rule uses $\gamma_y=0.25$; the new rule uses $\gamma_y=0$.
    Values are from the permanent-shock figure data at commits
    \texttt{fdc76c47} and \texttt{2f26bb3e}.}
\end{table}

\section{Why the inflation-only rule fixes the endpoint}

The revised calibration sets $\gamma_y=0$.
The endpoint Taylor rule is now
\begin{equation}
    \widehat r=\gamma_\pi\widehat\pi.
\end{equation}
Combining it with the Euler restriction $\widehat r=\widehat\pi$ yields
\begin{equation}
    (\gamma_\pi-1)\widehat\pi=0.
\end{equation}
Because $\gamma_\pi=1.50>1$, the unique endpoint is
\begin{equation}
    \boxed{\widehat\pi=0}.
\end{equation}
Permanent output gains no longer require permanent deflation.
Inflation can still move during the transition as marginal cost, demand, and expected
future conditions evolve.
The new simulations display modest positive inflation that declines toward target.
The productive-spending responses have not fully converged after 25 years, which is why
their inflation rates in Table~\ref{tab:comparison} are still above zero at that horizon.

\section{Interpretation}

The old result should be understood as a policy-rule artifact.
The central bank treated a permanent increase in sustainable output as a permanently
positive demand gap because the rule used $Y_t^d/\ssval{Y^d}$.
An output term could be reintroduced using an economically defined gap, such as output
relative to its natural level.
Such a gap would close after a permanent supply expansion, unlike the level deviation
used in equation~\eqref{eq:taylor-old}.

This investigation concerns the deterministic perfect-foresight experiments.
All recalibrated paths converge and contain finite values.
The extended model's pre-existing local rank-condition warning occurs under both values
of $\gamma_y$ and is therefore separate from the inflation mechanism derived here.

\end{document}
